%!TEX root = ../template.tex
%%%%%%%%%%%%%%%%%%%%%%%%%%%%%%%%%%%%%%%%%%%%%%%%%%%%%%%%%%%%%%%%%%%%
%% chapter7.tex
%% NOVA thesis document file
%%
%% Chapter with lots of dummy text
%%%%%%%%%%%%%%%%%%%%%%%%%%%%%%%%%%%%%%%%%%%%%%%%%%%%%%%%%%%%%%%%%%%%

\typeout{NT FILE chapter7.tex}%

\chapter{Unit Testing}
\label{cha:Unit Testing}

\section{Unit Tests: Trees and Attribute Annotations}
This chapter includes every unit test defined in the test module, with a brief purpose statement and a corresponding figure. 
Passing cases show fully annotated attribute values; failure-oriented cases indicate the precise failure point (type mismatch, division by zero, L-attributed 
violation, or out-of-range occurrence). Cycle-detection tests also include the attribute-dependency graph. 

\subsection{\texttt{test\_ag1\_simple\_synthesized} (\texttt{ag1}, tree \texttt{pt1})}
Simple synthesized value over multiplication: \(S\to E\), \(E\to E * F\), inner \(E\to F\); leaves \(F\to 3, F\to 2\). Expected \(v(S)=6\). 

\begin{verbatim}
Synthesized: v

S -> E {v(S) = v(E)}
E -> E * F {v(E0) = v(E1) * v(F)}
E -> F {v(E) = v(F)}
F -> 2 {v(F) = 2}
F -> 3 {v(F) = 3}

\end{verbatim}

\begin{tikzpicture}[
  level 1/.style={sibling distance=45mm, level distance=20mm},
  level 2/.style={sibling distance=30mm, level distance=20mm},
  level 3/.style={sibling distance=36mm, level distance=20mm}, 
  edge from parent/.style={draw, -latex},
  every node/.style={draw, ellipse, align=center}
]
\node {S\\$v=6$}
  child {node {E\\$v=6$}
    child {node {E\\$v=3$}
      child {node {F\\$v=3$}
        child {node {3}}
      }
    }
    child {node {*}}
    child {node {F\\$v=2$}
      child {node {2}}
    }
  };
\end{tikzpicture}

\vspace{1cm}

\subsection{\texttt{test\_parseTree} (Grammar \texttt{ag2}, word \texttt{3+2+9\textasciitilde})}
Parse \texttt{3+2+9\textasciitilde} from the CFG induced by \texttt{ag2}, convert to an AG tree, and evaluate attributes. Rules propagate a running sum via inherited \texttt{d} and synthesized \texttt{r}; result at root: \(v(S)=14\). 

\subsection{\texttt{test\_ag2\_simple\_synthesized\_and\_inherited} (\texttt{ag2}, tree \texttt{pt2})}
 Same structure as \texttt{test\_parseTree}; checks inherited \texttt{d} and synthesized \texttt{r} along the \texttt{X}-chain; \(v(S)=14\). 

\begin{verbatim}
Inherited: d
Synthesized: v, r

S -> E {v(S) = v(E)}
E -> F X {v(E) = r(X) ; d(X) = v(F)}
X -> + F X {r(X0) = r(X1) ; d(X1) = d(X0) + v(F)}
X -> ~ {r(X) = d(X)}
F -> 2 {v(F) = 2}
F -> 9 {v(F) = 9}
\end{verbatim}

\begin{tikzpicture}[
  level 1/.style={sibling distance=50mm, level distance=30mm},
  level 2/.style={sibling distance=40mm, level distance=30mm},
  level 3/.style={sibling distance=30mm, level distance=30mm},
  level 4/.style={sibling distance=20mm, level distance=30mm},
  edge from parent/.style={draw, -latex},
  every node/.style={draw, ellipse, align=center}
]
\node {S\\$v=14$}
  child {node {E\\$v=14$}
    child {node {F\\$v=3$}
      child {node {3}}
    }
    child {node {X\\$d=3$\\$r=14$}
      child {node {+}}
      child {node {F\\$v=2$}
        child {node {2}}
      }
      child {node {X\\$d=5$\\$r=14$}
        child {node {+}}
        child {node {F\\$v=9$}
          child {node {9}}
        }
        child {node {X\\$d=14$\\$r=14$}
          child {node {\textasciitilde}}
        }
      }
    }
  };
\end{tikzpicture}


\subsection{\texttt{test\_ag3\_simpler\_synthesized\_and\_inherited} (\texttt{ag3}, tree \texttt{pt3})}
 Minimal inherited/synthesized flow: \(S\to E\ \{d(E)=5,\ v(S)=v(E)\}\), \(E\to \sim\ \{v(E)=d(E)+1\}\). Expected \(v(S)=6\). 

\begin{verbatim}
Inherited: d
Synthesized: v

S -> E {v(S) = v(E) ; d(E) = 5}
E -> ~ {v(E) = d(E) + 1}
\end{verbatim}

\begin{tikzpicture}[
  level 1/.style={sibling distance=35mm, level distance=30mm},
  edge from parent/.style={draw, -latex},
  every node/.style={draw, ellipse, align=center}
]
\node {S\\$v=6$}
  child {node {E\\$d=5$\\$v=6$}
    child {node {\textasciitilde}}
  };
\end{tikzpicture}

\vspace{1cm}


\subsection{\texttt{test\_ag4\_type\_mismatch} (\texttt{ag4}, tree \texttt{pt4})}
Deliberate type error: \(S\to A\ B\ \{\,b(S)=b(A)+b(B)\,\}\) with \(A\to x\{b(A)=\mathrm{T}\}\), \(B\to y\{b(B)=\mathrm{F}\}\). The sum \(+\) is invalid on booleans; failure at the root. 

\begin{verbatim}
Synthesized: b

S -> A B { b(S) = b(A) + b(B) }
A -> x { b(A) = T }
B -> y { b(B) = F }
\end{verbatim}


\begin{tikzpicture}[
  level 1/.style={sibling distance=44mm},
  level 2/.style={sibling distance=26mm},
  edge from parent/.style={draw, -latex},
  every node/.style={draw, ellipse, align=center}
]
\node {S\\$b=\textnormal{type error}$}
  child {node {A\\$b=\mathrm{T}$}
    child {node {x}}
  }
  child {node {B\\$b=\mathrm{F}$}
    child {node {y}}
  };
\end{tikzpicture}




\subsection{\texttt{test\_ag5\_div\_zero} (\texttt{ag5}, tree \texttt{pt5})}
 Division by zero: \(E\to E / F\) with left \(F\to 1\) and right \(F\to 0\). Failure occurs evaluating \(v(E_0)=v(E_1)/v(F)\). 


\begin{verbatim}
Synthesized: v

S -> E {v(S) = v(E)}
E -> E / F {v(E0) = v(E1) / v(F)}
E -> F {v(E) = v(F)}
F -> 0 {v(F) = 0}
F -> 1 {v(F) = 1}
\end{verbatim}

\begin{tikzpicture}[
  level 1/.style={sibling distance=48mm, level distance=20mm},
  level 2/.style={sibling distance=30mm, level distance=20mm},
  level 3/.style={sibling distance=22mm, level distance=20mm},
  edge from parent/.style={draw, -latex},
  every node/.style={draw, ellipse, align=center}
]
\node {S\\$v=\textnormal{div-by-zero}$}
  child {node {E\\$v=\textnormal{div-by-zero}$}
    child {node {E\\$v=1$}
      child {node {F\\$v=1$}
        child {node {1}}
      }
    }
    child {node {/}}
    child {node {F\\$v=0$}
      child {node {0}}
    }
  };
\end{tikzpicture}



\subsection{\texttt{test\_ag6\_lattr\_violation} (\texttt{ag6}, tree \texttt{pt6})}
 L-attributed violation at runtime: \(S\to A\ B\ \{d(A)=v(B)\}\) requires reading \(v(B)\) before \(B\) is evaluated (left-to-right discipline). Failure occurs when setting \(d(A)\). 

\begin{verbatim}
Inherited: d
Synthesized: v

S -> A B { d(A) = v(B) ; v(S) = v(A) + v(B) }
A -> a { v(A) = d(A) }
B -> a { v(B) = 1 }
\end{verbatim}


\begin{tikzpicture}[
  level 1/.style={sibling distance=44mm, level distance=20mm},
  level 2/.style={sibling distance=26mm, level distance=20mm},
  edge from parent/.style={draw, -latex, level distance=20mm},
  every node/.style={draw, ellipse, align=center}
]
\node {S\\$v=\textnormal{not computed}$}
  child {node {A\\$d=\textnormal{needs }v(B)$\\$v=\textnormal{blocked}$}
    child {node {a}}
  }
  child {node {B\\$v=1$}
    child {node {a}}
  };
\end{tikzpicture}



\subsection{\texttt{test\_ag7\_default\_index} (\texttt{ag7}, tree \texttt{pt7})}
 Default vs.\ explicit child indices on \(S\to C\ C\); inherited \(d(C_1)=1\), \(d(C_2)=3\); \(v(S)=v(C_1)+v(C_2)=4\). 


\begin{verbatim}
Inherited: d
Synthesized: v

S -> C C { d(C) = 1 ; d(C2) = 3 ; v(S) = v(C1) + v(C2) }
C -> n { v(C) = d(C) }
C -> m { v(C) = d(C) }
\end{verbatim}

\begin{tikzpicture}[
  level 1/.style={sibling distance=44mm, level distance=20mm},
  level 2/.style={sibling distance=26mm, level distance=20mm},
  edge from parent/.style={draw, -latex},
  every node/.style={draw, ellipse, align=center}
]
\node {S\\$v=4$}
  child {node {C$_1$\\$d=1$\\$v=1$}
    child {node {n}}
  }
  child {node {C$_2$\\$d=3$\\$v=3$}
    child {node {m}}
  };
\end{tikzpicture}


\subsection{\texttt{test\_ag8\_boolean\_flags} (\texttt{ag8}, tree \texttt{pt8})}
 Boolean synthesized attribute \texttt{b} with \(A\to x\{b(A)=\mathrm{T}\}\), \(B\to y\{b(B)=\mathrm{F}\}\), and \(S\to A\,B\) with equality check; \(b(S)=\mathrm{F}\). 


\begin{verbatim}
Synthesized: b

S -> A B { b(S) = b(A) = b(B) }
A -> x { b(A) = T }
A -> y { b(A) = F }
B -> x { b(B) = T }
B -> y { b(B) = F }
\end{verbatim}

\begin{tikzpicture}[
  level 1/.style={sibling distance=44mm},
  level 2/.style={sibling distance=26mm},
  edge from parent/.style={draw, -latex},
  every node/.style={draw, ellipse, align=center}
]
\node {S\\$b=\mathrm{F}$}
  child {node {A\\$b=\mathrm{T}$}
    child {node {x}}
  }
  child {node {B\\$b=\mathrm{F}$}
    child {node {y}}
  };
\end{tikzpicture}


% ------------------------------------------------------------
\subsection{\texttt{test\_ag9\_list\_length} (\texttt{ag9}, tree \texttt{pt9}=\texttt{make\_list 30})}
 List length via synthesized \texttt{v}. For \(n=30\) the root has \(v=30\). Below is the same left-nested pattern shown for \(n=4\) (illustrative); the test uses \(n=30\). 

\begin{verbatim}
Synthesized: v

L -> L , I { v(L0) = v(L1) + v(I) }
L -> I { v(L) = v(I) }
I -> a { v(I) = 1 }
\end{verbatim}

\begin{tikzpicture}[
  level 1/.style={sibling distance=40mm, level distance=20mm},
  level 2/.style={sibling distance=35mm, level distance=20mm},
  level 3/.style={sibling distance=28mm, level distance=20mm},
  level 4/.style={sibling distance=22mm, level distance=20mm},
  edge from parent/.style={draw, -latex},
  every node/.style={draw, ellipse, align=center}
]
\node {L\\$v=4$}
  child {node {L\\$v=3$}
    child {node {L\\$v=2$}
      child {node {L\\$v=1$}
        child {node {I\\$v=1$}
          child {node {a}}
        }
      }
      child {node {,}}
      child {node {I\\$v=1$}
        child {node {a}}
      }
    }
    child {node {,}}
    child {node {I\\$v=1$}
      child {node {a}}
    }
  }
  child {node {,}}
  child {node {I\\$v=1$}
    child {node {a}}
  };
\end{tikzpicture}




% ------------------------------------------------------------
\subsection{\texttt{test\_ag10\_inherited\_default} (\texttt{ag10}, tree \texttt{pt10})}
 Default inherited values \(d(C)=0,d(F)=0\), then \(v(C)=1\), \(v(F)=2\); hence \(v(S)=3\). 

 
\begin{verbatim}
Inherited: d
Synthesized: v

S -> C F { d(C) = 0; d(F) = 0 ; v(S) = v(C) + v(F) }
C -> n { v(C) = d(C) + 1 }
F -> m { v(F) = d(F) + 2 }
\end{verbatim}


\begin{tikzpicture}[
  level 1/.style={sibling distance=46mm, level distance=20mm},
  level 2/.style={sibling distance=26mm, level distance=20mm},
  edge from parent/.style={draw, -latex},
  every node/.style={draw, ellipse, align=center}
]
\node {S\\$v=3$}
  child {node {C\\$d=0$\\$v=1$}
    child {node {n}}
  }
  child {node {F\\$d=0$\\$v=2$}
    child {node {m}}
  };
\end{tikzpicture}


% ------------------------------------------------------------
\subsection{\texttt{test\_ag11\_out\_of\_range\_ref} (\texttt{ag11}, tree \texttt{pt11})}
 Out-of-range child reference: \(S\to A\ \{v(S)=v(A_2)\}\) but only one \(A\) child exists. Failure at the root while resolving \(A_2\). 

\begin{verbatim}
Synthesized: v

S -> A { v(S) = v(A2) }
A -> a { v(A) = 1 }
\end{verbatim}

\begin{tikzpicture}[
  level 1/.style={sibling distance=40mm, level distance=20mm},
  level 2/.style={sibling distance=24mm, level distance=20mm},
  edge from parent/.style={draw, -latex},
  every node/.style={draw, ellipse, align=center}
]
\node {S\\$v=\textnormal{ref A\_2 out of range}$}
  child {node {A\\$v=1$}
    child {node {a}}
  };
\end{tikzpicture}


\subsection{\texttt{test\_has\_cycles\_ok} (\texttt{ok\_ag})}
 Acyclic grammar \(\;S\to F\{v(S)=v(F)\},\ F\to a\{v(F)=1\}\;\). No dependency cycle; a concrete tree for the terminal word \texttt{a} is shown. (The unit test itself only checks the cycle predicate.) 

\begin{verbatim}
Synthesized: v

S -> F { v(S) = v(F) }
F -> a { v(F) = 1 }
\end{verbatim}


\begin{tikzpicture}[
  level 1/.style={sibling distance=36mm, level distance=20mm},
  level 2/.style={sibling distance=24mm, level distance=20mm},
  edge from parent/.style={draw, -latex},
  every node/.style={draw, ellipse, align=center}
]
\node {S\\$v=1$}
  child {node {F\\$v=1$}
    child {node {a}}
  };
\end{tikzpicture}




% ------------------------------------------------------------
\subsection{\texttt{test\_has\_cycles\_detected} (\texttt{cyc\_ag})}
 Deliberate attribute cycle: \(S\to F\{v(S_0)=v(F_0)\},\ F\to S\{v(F_0)=v(S_0)\}\). No finite parse tree exists; the test inspects the dependency graph (cycle between \(v(S_0)\) and \(v(F_0)\)). 

\begin{verbatim}
Synthesized: v

S -> F { v(S0) = v(F0) }
F -> S { v(F0) = v(S0) }
\end{verbatim}

\begin{center}
\begin{tikzpicture}[
  node distance=40mm,
  edge from parent/.style={draw, -latex},
  every node/.style={draw, ellipse, align=center, minimum width=22mm, minimum height=8mm}
]
\node (s) {$v(S_0)$};
\node (f) [right=of s] {$v(F_0)$};
\draw[-latex] (s) -- (f);
\draw[-latex] (f) -- (s);
\end{tikzpicture}
\end{center}


% ------------------------------------------------------------
\subsection{\texttt{test\_accept\_words}}
 Checks CFG acceptance for the word \texttt{3\*2} derived from \texttt{ag1}. The test does not evaluate attributes, but a consistent annotated AG tree would have \(v(S)=6\). 

\begin{tikzpicture}[
  level 1/.style={sibling distance=46mm, level distance=20mm},
  level 2/.style={sibling distance=30mm, level distance=20mm},
  level 3/.style={sibling distance=22mm, level distance=20mm},
  edge from parent/.style={draw, -latex},
  every node/.style={draw, ellipse, align=center}
]
\node {S\\$v=6$}
  child {node {E\\$v=6$}
    child {node {E\\$v=3$}
      child {node {F\\$v=3$}
        child {node {3}}
      }
    }
    child {node {*}}
    child {node {F\\$v=2$}
      child {node {2}}
    }
  };
\end{tikzpicture}

% ------------------------------------------------------------
\subsection{\texttt{test\_validate}}
 Runs static validation on the extended arithmetic grammar \texttt{ag\_expr\_ext}. A representative evaluated tree for the fixture \texttt{pt\_expr\_ext} (expression \(3 + 2 \times (9-5) / 2 - 1\)) is shown in Section~\ref{sec:expr-ext-tree}. (Validation itself performs no tree evaluation.) 


% ------------------------------------------------------------
\subsection{\texttt{test\_ag\_expr\_ext} (\texttt{ag\_expr\_ext}, tree \texttt{pt\_expr\_ext})}
\label{sec:expr-ext-tree}
 Extended arithmetic with attributes value \texttt{v} and pretty string \texttt{s}. Root result: \(v(S)=6\), \(s(S)=\)\texttt{3+2*(9-5)/2-1}. (Height \texttt{h} is computed by the grammar but omitted here for brevity.) 

 \begin{verbatim}
Synthesized: v, h, s

S -> E { v(S) = v(E) ; s(S) = s(E) ; h(S) = h(E) }

E -> E + T { v(E0) = v(E1) + v(T) ; s(E0) = s(E1) + '+' + s(T) ;
             h(E0) = h(E1) + h(T) + 1 }
E -> E - T { v(E0) = v(E1) - v(T) ; s(E0) = s(E1) + '-' + s(T) ; 
             h(E0) = h(E1) + h(T) + 1 }
E -> T     { v(E)  = v(T) ; s(E)  = s(T) ; h(E)  = h(T) }

T -> T * F { v(T0) = v(T1) * v(F) ; s(T0) = s(T1) + '*' + s(F) ;
             h(T0) = h(T1) + h(F) + 1 }
T -> T / F { v(T0) = v(T1) / v(F) ; s(T0) = s(T1) + '/' + s(F) ;
             h(T0) = h(T1) + h(F) + 1 }
T -> F     { v(T)  = v(F) ; s(T)  = s(F) ; h(T)  = h(F) }

F -> ( E ) { v(F) = v(E) ; s(F) = '(' + s(E) + ')' ; h(F) = h(E) + 1 }
F -> N     { v(F) = v(N) ; s(F) = s(N) ; h(F) = h(N) }

N -> 0 { v(N) = 0 ; s(N) = '0' ; h(N) = 1 }
N -> 1 { v(N) = 1 ; s(N) = '1' ; h(N) = 1 }
N -> 2 { v(N) = 2 ; s(N) = '2' ; h(N) = 1 }
N -> 3 { v(N) = 3 ; s(N) = '3' ; h(N) = 1 }
N -> 4 { v(N) = 4 ; s(N) = '4' ; h(N) = 1 }
N -> 5 { v(N) = 5 ; s(N) = '5' ; h(N) = 1 }
N -> 6 { v(N) = 6 ; s(N) = '6' ; h(N) = 1 }
N -> 7 { v(N) = 7 ; s(N) = '7' ; h(N) = 1 }
N -> 8 { v(N) = 8 ; s(N) = '8' ; h(N) = 1 }
N -> 9 { v(N) = 9 ; s(N) = '9' ; h(N) = 1 }
\end{verbatim}

\begin{tikzpicture}[
  level 1/.style={sibling distance=62mm, level distance=40mm},
  level 2/.style={sibling distance=46mm, level distance=20mm},
  level 3/.style={sibling distance=34mm, level distance=20mm},
  level 4/.style={sibling distance=28mm, level distance=20mm},
  level 5/.style={sibling distance=24mm, level distance=20mm},
  edge from parent/.style={draw, -latex},
  every node/.style={draw, ellipse, align=center}
]
\node {S\\$v=6$\\$s=\texttt{3+2*(9-5)/2-1}$}
  child {node {E\\$v=6$}
    child {node {E\\$v=7$}
      child {node {E\\$v=3$}
        child {node {T\\$v=3$}
          child {node {F\\$v=3$}
            child {node {N\\$v=3$}
              child {node {3}}
            }
          }
        }
      }
      child {node {+}}
      child {node {T\\$v=4$}
        child {node {T\\$v=8$}
          child {node {T\\$v=2$}
            child {node {F\\$v=2$}
              child {node {N\\$v=2$}
                child {node {2}}
              }
            }
          }
          child {node {*}}
          child {node {F\\$v=4$}
            child {node {(}}
            child {node {E\\$v=4$}
              child {node {E\\$v=9$}
                child {node {T\\$v=9$}
                  child {node {F\\$v=9$}
                    child {node {N\\$v=9$}
                      child {node {9}}
                    }
                  }
                }
              }
              child {node {-}}
              child {node {T\\$v=5$}
                child {node {F\\$v=5$}
                  child {node {N\\$v=5$}
                    child {node {5}}
                  }
                }
              }
            }
            child {node {)}}
          }
        }
        child {node {/}}
        child {node {F\\$v=2$}
          child {node {N\\$v=2$}
            child {node {2}}
          }
        }
      }
    }
    child {node {-}}
    child {node {T\\$v=1$}
      child {node {F\\$v=1$}
        child {node {N\\$v=1$}
          child {node {1}}
        }
      }
    }
  };
\end{tikzpicture}